\chapter{Orchestrators, Job-Local Storage and Data Movement}
\label{ch:hpc-storage}

\Cref{ch:backends} placed tasks in Docker containers and Slurm jobs. This
chapter completes the infrastructure side of the course (Lectures~38--40 and
42--44). It covers HPC containers with Apptainer, container orchestrators,
remote and cloud hosts, job-local BeeOND file systems, pre-staging of inputs
before submission, and the data movers that carry \texttt{workflow://}
references between hosts. Most of these services are not on a laptop. Each
example therefore checks what DAGonStar generates or decides, and runs the
rest under portable emulation or against a local stand-in.

\section{Apptainer}

\texttt{TaskType.APPTAINER} runs a command inside a Singularity Image Format
(SIF) file. If the image is a registry reference (the default is
\texttt{docker://ubuntu:20.04}), the task builds a \texttt{.sif} with
\texttt{apptainer build}, retrying with \texttt{sudo} if needed, and reuses an
existing file afterwards. The command runs with \texttt{apptainer exec}, a
writable overlay (\texttt{overlay\_size} 1024\,MiB by default), the scratch
directory bound at the same path, and \texttt{--pwd} set to the task folder.

\lstinputlisting[caption={Lecture 38: an Apptainer task (\texttt{examples/38-apptainer/apptainer\_task.py}).},label={lst:apptainer}]{../../examples/38-apptainer/apptainer_task.py}

\lstinputlisting[style=bash,caption={Recorded output of Lecture 38.}]{../../examples/38-apptainer/output/apptainer_task.txt}

\section{Kubernetes and Nomad}

A \texttt{KUBERNETES} task creates a pod that runs \texttt{sleep infinity},
executes the launcher in the container \texttt{main} with
\texttt{/bin/bash -c}, and deletes the pod afterwards. A \texttt{NOMAD} task
submits a batch job with the Docker driver and a restart policy of zero
attempts, so a failure is reported, not hidden. Its log is fetched and the job
is purged. Lecture~39 replaces both clients with local fakes and checks the
objects DAGonStar sends.

\lstinputlisting[style=bash,caption={Recorded output of Lecture 39 (\texttt{examples/39-orchestrators/kubernetes\_nomad.py}).}]{../../examples/39-orchestrators/output/kubernetes_nomad.txt}

\section{Remote Hosts and Clouds}

Passing \texttt{ip} and \texttt{ssh\_username} turns a batch, Slurm or Docker
task into its remote variant (\Cref{sec:factory-pattern-depth}). A connection
is opened only when both are given. Every argument of a remote command is
quoted, so a folder name with spaces or a semicolon cannot become a second
command. \texttt{CLOUD} tasks use Apache Libcloud, which is an optional extra.
Without it the task fails with a message that names the install command.

\lstinputlisting[style=bash,caption={Recorded output of Lecture 40 (\texttt{examples/40-remote/ssh\_and\_cloud.py}).}]{../../examples/40-remote/output/ssh_and_cloud.txt}

\section{BeeOND: Job-Local File Systems}

BeeOND builds a parallel file system from the local disks of a job's nodes
for the job's duration. A Slurm task requests it with \texttt{beeond="task"}
for a private volume. With \texttt{beeond="workflow:<name>"} it uses a volume
created once with \texttt{wf.create\_beeond\_volume}. The task finds its volume
in the \texttt{DAGON\_BEEOND} environment variable. With \texttt{persist}
(the default), the contents of a task volume are copied into the task folder
under \texttt{beeond/}, so consumers reference them as
\texttt{workflow:///<task>/beeond/<file>}.

\lstinputlisting[caption={Lecture 42: task and workflow volumes (\texttt{examples/42-beeond/beeond.py}).},label={lst:beeond}]{../../examples/42-beeond/beeond.py}

\lstinputlisting[style=bash,caption={Recorded output of Lecture 42.}]{../../examples/42-beeond/output/beeond.txt}

\begin{cautionbox}[A shared volume carries no references]
Tasks that share a workflow volume exchange files through
\texttt{\$DAGON\_BEEOND}, not through \texttt{workflow://} references. The
engine cannot infer an edge from them, so the order must be declared with
\texttt{add\_dependency\_to}, as Lecture~42 does for \texttt{roms} and
\texttt{stage}.
\end{cautionbox}

\section{Pre-Staging Before Submission}

On a busy cluster a job may wait hours in the queue and then fail in its
first second because an input is missing. With \texttt{pre\_stage=True}
DAGonStar stages and verifies a Slurm task's inputs \emph{before}
\texttt{sbatch}. The phases are \texttt{PREPARING}, \texttt{STAGING},
\texttt{STAGED}, \texttt{SUBMITTED}, then \texttt{COMPLETED} or
\texttt{FAILED}. A completion marker (\texttt{stage-in-complete.json}) must
exist before submission. A failed stage-in never reaches the scheduler.
Pre-staging is refused at \texttt{add\_task} for backends that cannot honour
it, and for tasks whose inputs live on a BeeOND volume that exists only inside
the job.

\lstinputlisting[caption={Lecture 43: pre-staging against a local stand-in for \texttt{sbatch} (\texttt{examples/43-prestage/prestage.py}).},label={lst:prestage}]{../../examples/43-prestage/prestage.py}

\lstinputlisting[style=bash,caption={Recorded output of Lecture 43.}]{../../examples/43-prestage/output/prestage.txt}

\section{Data Movers}

When a consumer and its producer run on the same host, the stager copies
(\texttt{cp -r}) or links (\texttt{ln -sf}) the referenced file. Across hosts
it tries GridFTP, then SCP, then FTP. Globus endpoints are used when both
hosts declare one. The mover mode (\texttt{NORMAL}, \texttt{PARALLEL},
\texttt{SLURM}) decides whether transfers run in sequence, in parallel or as
a Slurm job.

\lstinputlisting[style=bash,caption={Recorded output of Lecture 44 (\texttt{examples/44-movers/data\_movers.py}).}]{../../examples/44-movers/output/data_movers.txt}

\section{Chapter Summary}

\begin{itemize}[leftmargin=1.4em]
  \item Apptainer, Kubernetes and Nomad change where a launcher runs, not the
    graph around it.
  \item Remote commands are quoted, and cloud support is an optional extra.
  \item BeeOND volumes are fast and temporary. Persist what consumers need,
    and order shared-volume tasks by hand.
  \item Pre-staging moves input failures from inside the job to before
    submission, where they cost no core-hours.
  \item The stager picks the cheapest transport that the two hosts support.
\end{itemize}

\begin{exercise}
A WRF run of 128 cores waits on average three hours in the queue, and one run
in twenty fails at start-up because a boundary file is missing. Estimate the
core-hours that pre-staging saves over a hundred runs, and say what it does
not save.
\end{exercise}
